\begin{titlepage}
  \begin{center}
    % Kopf der Seite
    \begin{tabular}{cl}
    \parbox[b]{10.2cm}{
      % \textsf würde das Aussehen der ersten Seite ruinieren, 
      % wer will, soll das selbst außen rum machen...
      \textbf{Hochschule RheinMain}\\
      Fachbereich Design Informatik Medien\\
      Studiengang \welcherstudiengang\ }
& 
\parbox[t]{6cm}{
  \includegraphics[width=0.27\textwidth]{figures/hsrmlogo.pdf}
  }
    \end{tabular}
    \vfill    
    {\LARGE \welchethesis-Arbeit} \\[0.5cm]
    {\large zur Erlangung des akademischen Grades} \\[5mm]
    {\large \welchethesis\ \thesisofwas} \\[5mm]
    \rule{\textwidth}{1pt}\\[0.5cm]
    {\begin{spacing}{1.15} \huge \bfseries \titel \\ \end{spacing}}
    \rule{\textwidth}{1pt}    
    \vfill    
    \begin{tabular}{ll} % Mitte der Seite
      Vorgelegt von & \autor \\
      am & \datum \\
      Referent & \referent \\
      Korreferent & \korreferent
    \end{tabular}    
    \vfill
  \end{center}
\end{titlepage}
\cleardoublepage

\thispagestyle{empty}
\section*{Eigenständigkeitserklärung\newline (gem. ABPO, Ziff. 4.1.5.4 bzw. RPO, §19.6)}

Ich bin verantwortlich für die Qualität des Inhalts dieser Arbeit, den ich mit geeigneten wissenschaftlichen Quellen belegt bzw. gestützt habe. Die aus fremden Quellen direkt oder indirekt übernommenen Texte, Gedankengänge, Konzepte, Grafiken, technischen Inhalte und Ähnliches in meinen Ausführungen habe ich eindeutig gekennzeichnet und mit vollständigen Verweisen auf die jeweilige Quelle versehen. Alle weiteren Inhalte der Arbeit (Textteile, Abbildungen, Tabellen etc.) ohne entsprechende Verweise stammen im urheberrechtlichen Sinn von mir.

Ich versichere außerdem, dass ich die \welchethesis-Arbeit selbständig verfasst und keine anderen als die angegebenen Quellen und Hilfsmittel benutzt sowie keine unerlaubte fremde Hilfe in Anspruch genommen habe.

%%% Anfang rote Hinweisbox zur KI-Verwendung
%%% vor Abgabe durch vereinbarte Formulierung ersetzen, Vorschläge s.u.
\begin{center}
  {\color{red}\fbox{\parbox{0.8\linewidth}{\it 
  Eine zu verwendende Formulierung zur Verwendung generativer KIs ist mit dem Referenten abzustimmen. Mögliche Formulierungen sind in der \LaTeX-Vorlage auskommentiert enthalten. Auch hiervon abweichende Regelungen sind in Absprache mit dem Referenten möglich. Bitte hier die vereinbarte Formulierung eintragen und diese rote Hinweisbox entfernen.
  }}}
\end{center}
%%% ENDE rote Hinweisbox zur KI-Verwendung

%%% Bei Verbot KI-generierter Inhalte %%%
% Ich versichere, dass ich die hier vorliegende Arbeit vollständig eigenständig erarbeitet,
% also keine KI-generierten Inhalte verwendet habe.

%%% Bei Erlaubnis KI-generierter Inhalte %%%
% Ich versichere, dass ich KI-Tools lediglich als Hilfsmittel verwendet habe und in der 
% vorliegenden Arbeit mein gestalterischer Einfluss überwiegt. Ich verantworte die Übernahme 
% jeglicher von mir verwendeter KI-generierter Inhalte vollumfänglich selbst. 

%%% (optional, falls Kennzeichnungspflicht) %%%
% Unter ``Übersicht verwendeter Hilfsmittel'' habe ich im Anhang der Arbeit sämtliche 
% KI-generierten Inhalte einzeln aufgeführt und die jeweils verwendete KI-Anwendung benannt. 
  
%%% (optional, falls Prompts abzugeben) %%%
% Die jeweiligen Benutzungseingaben (z.B. Prompts) sind der Arbeit beigefügt.


Die vorliegende Arbeit wurde bisher weder im In- noch im Ausland in gleicher oder ähnlicher Form einer anderen Prüfungsbehörde vorgelegt. Mir ist bekannt, dass ein Verstoß gegen die genannten Punkte als Täuschungsversuch gelten und prüfungsrechtliche Konsequenzen haben kann. Insbesondere kann es dazu führen, dass die Leistung nicht bestanden ist und dass bei mehrfachem oder schwerwiegendem Täuschungsversuch eine Exmatrikulation droht.

\vfill
\noindent\begin{tabular}{p{0.37\textwidth}p{0.56\textwidth}}
\ort, \datum  & \rule{0.56\textwidth}{0.5pt}\\
              & \makebox[1cm]{\ } \autor
\end{tabular}

\vfill\mbox{}
\cleardoublepage

\section*{Erklärung zur Verwendung der \welchethesis-Arbeit}

Hiermit erkläre ich mein Einverständnis mit den im folgenden 
aufgeführten Verbreitungsformen dieser Abschlussarbeit:
\bigskip

\noindent\begin{tabular}{|p{0.82\textwidth}|c|c|}
  \hline
  \textbf{Verbreitungsform} & \makebox[0.035\textwidth]{\textbf{Ja}} 
                            & \makebox[0.05\textwidth]{\textbf{Nein}} \\\hline
  Veröffentlichung des Titels der Arbeit im Internet  
                                           & $\times$ & \\\hline
  Veröffentlichung der Arbeit im Internet             
                                           &  & $\times$ \\\hline
\end{tabular}

\vspace{2cm}
\noindent\begin{tabular}{p{0.37\textwidth}p{0.56\textwidth}}
\ort, \datum  & \rule{0.56\textwidth}{0.5pt}\\
              & \makebox[1cm]{\ } \autor
\end{tabular}
\cleardoublepage



